\section{Limitations}
\label{sec:limitations}

Several protocol choices narrow how far these results generalize. Sample size is the first constraint. The prompt-injection test set has only $n=116$ rows, so accuracy and F1 intervals are wide and small ranking differences are unstable. JailbreakBench uses 100 harmful and 100 benign prompts; HarmBench and other harmful-only sets contribute miss rate or detection rate without a paired false-positive rate. We do not headline a single winner from the smallest benchmarks.

Threshold fitting follows pre-registered rules that sometimes confirm the operating point at 0.5. When a labeled development pool has fewer than 200 rows, tuned thresholds are not applied and metrics are reported at 0.5. JailbreakBench falls in that bucket, so all decision-model and stack numbers on that set use the fixed 0.5 cutoffs on test. That choice avoids overfitting a threshold to a tiny split but may understate what development tuning could achieve elsewhere.

WildGuardMix shows a second threshold issue. The benign development pool for prompt-harm tuning is very small. The rule that tunes score thresholds to maximize harmful-class F1 subject to at most 5\% benign false-positive rate on development data can move the cutoff away from 0.5 and lower test F1 relative to the 0.5 baseline. On this run Jev prompt-harm F1 is 0.4940 tuned versus 0.5904 at 0.5, even though prompt-harm AUPRC stays high. That gap is a limitation of the tuning budget and pool size, not proof that 0.5 is always better on WildGuardMix in production.

Input length is handled differently by provider. Jev records provider context overflow as an error and keeps the row in the metric table. Clef and Clef-flash on Workers AI silently shorten an over-long \texttt{state} and still return an answer; the scorer marks those rows \texttt{truncated} when token counts exceed the hosted limits (24{,}576 tokens for Clef-flash and 64{,}000 for Clef). We do not separately measure whether truncation changes scores on the truncated tail of long prompts, and we do not rerun Clef-flash with self-hosted weights at a different context window.

The study scores prompts only. We do not generate new attacks and we do not measure whether a downstream chat model would comply with a harmful instruction if the screen allowed it. HarmBench attack-success rates against a target model are out of scope. Miss rate here means the screen allowed the text; it is a filter outcome, not an end-to-end safety outcome.

Guardrail outputs are ordinal confidence bands, not calibrated probabilities. We do not compute expected calibration error or Brier scores for Guardrails, and we do not claim that a band is ``better calibrated'' than Jev \texttt{P(yes)}. Comparisons across arms mix probabilistic decision-model scores with vendor-specific band rules at low, medium, and high strength.

Guardrails run in \texttt{us-east-2} on a pinned version. Costs use 2026-10-19 list prices. Latency uses stored \texttt{latency\_ms} from the original call; cache replays keep that value rather than zeroing it.

Dataset label noise and possible training-data overlap on public benchmarks further limit causal claims about unseen deployment traffic.
